\chapter*{前言}
\BookTopBookmark{前言}{book-preface}
\markboth{前言}{}

\section*{本书目标}

机器学习把数据、模型和计算连接起来，但理解一种方法，不能只记住模型名称或训练步骤。
还需要知道模型能够表示什么规律，学习结论依赖哪些数据与假设，什么证据足以支持评价，
以及方法进入具体应用和运行系统后会受到哪些约束。本书沿着这些联系组织内容，希望读者
不仅会使用算法，也能解释一个结果为何成立、何时可信，以及怎样把它放入完整的技术过程。

本书面向已经具备基本微积分、线性代数、概率统计和编程基础，希望系统理解机器学习的
读者。正文会在需要时回顾具体工具，但不以这些基础课程的完整讲解为目标。阅读后，读者
应能从任务和证据出发比较模型与学习方式，并把理论条件、实验评价、应用要求和系统实现
放在同一条推理链上考察。

\section*{本书内容}

全书分为五卷。第一卷“基础、理论和可信性”建立任务、模型、学习与评价的共同语言，
说明学习保证依赖的条件，并讨论可靠性、鲁棒性、公平性、隐私等可信要求应由什么证据
支持。第二卷“模型”回答规律可以怎样表示，依次考察经典模型、神经网络模型和概率图
模型的结构、推断与学习。第三卷“范式”转向经验和知识的获得方式，讨论交互反馈、有限
标注以及任务变化下如何学习、复用和迁移知识。
前面三卷提供分析工具和方法选择，后两卷把这些选择放入实际约束。第四卷“应用”围绕
自然语言处理、图像处理、推荐与搜索，说明任务目标、数据形态、模型和评价方式如何相互
制约。第五卷“系统”连接数据准备、训练计算、部署服务与运行维护，考察方法如何成为可
持续运行的机器学习系统。这样的次序并不表示建模、应用和系统彼此割裂：应用需求决定
评价口径，系统条件限制可用方法，而运行中的证据又会促使模型和学习过程重新调整。

\section*{参考资料}

下面按内容列出可与本书配合阅读的教材。每本书的侧重点不同，读者可以根据正在阅读的
卷和需要补充的问题选择查阅，不必按列表顺序通读。

\paragraph*{综合与基础}
\begin{itemize}
  \item Mitchell 的 \emph{Machine Learning} 从概念学习、归纳偏置和代表性算法出发，
    适合查阅机器学习的基本问题和早期统一框架。
  \item Bishop 的 \emph{Pattern Recognition and Machine Learning} 将概率建模、
    决策理论、核方法和近似推断联系起来，适合查阅概率视角下的模型与推断方法。
  \item 周志华的《机器学习》覆盖常用模型、学习理论和实践问题，适合按任务比较方法的
    基本假设与用途。
  \item 李航的《统计学习方法（第 2 版）》以统计机器学习为定位，适合按模型查阅定义、
    算法和数学推导。
  \item Deisenroth、Faisal 和 Ong 的 \emph{Mathematics for Machine Learning}
    适合补充线性代数、解析几何、矩阵微积分、概率与优化等数学工具。
\end{itemize}

\paragraph*{统计学习与学习理论}
\begin{itemize}
  \item Hastie、Tibshirani 和 Friedman 的 \emph{The Elements of Statistical
    Learning} 适合查阅回归、分类、正则化、模型选择和统计学习方法。
  \item Shalev-Shwartz 和 Ben-David 的 \emph{Understanding Machine Learning:
    From Theory to Algorithms} 适合查阅PAC框架、复杂度、泛化界与优化基础。
  \item Mohri、Rostamizadeh 和 Talwalkar 的 \emph{Foundations of Machine Learning}
    适合查阅学习界、核方法、在线学习及其理论保证。
  \item James、Witten、Hastie 和 Tibshirani 的
    \emph{An Introduction to Statistical Learning} 适合结合R实例理解回归、分类、
    重采样、正则化、树模型和无监督学习。
\end{itemize}

\paragraph*{概率模型与图模型}
\begin{itemize}
  \item Murphy 的 \emph{Machine Learning: A Probabilistic Perspective} 以概率视角
    组织监督学习、无监督学习、潜变量模型和图模型，适合查阅模型间的联系。
  \item Murphy 的 \emph{Probabilistic Machine Learning: An Introduction} 适合查阅现代
    概率机器学习的统一记号、计算方法和模型实例。
  \item Murphy 的 \emph{Probabilistic Machine Learning: Advanced Topics}
    适合查阅深度生成模型、状态空间模型、贝叶斯优化和相关近似推断方法。
  \item Koller 和 Friedman 的 \emph{Probabilistic Graphical Models: Principles and
    Techniques} 适合系统查阅图模型的表示、精确与近似推断以及参数和结构学习。
  \item Rasmussen 和 Williams 的 \emph{Gaussian Processes for Machine Learning}
    适合深入理解高斯过程、核函数、贝叶斯预测和超参数学习。
\end{itemize}

\paragraph*{神经网络与深度学习}
\begin{itemize}
  \item Goodfellow、Bengio 和 Courville 的 \emph{Deep Learning} 适合查阅神经网络
    表示、优化、正则化和深度生成模型。
  \item Bishop 和 Bishop 的 \emph{Deep Learning: Foundations and Concepts} 从概率
    基础连接常用架构与训练方法，适合查阅相应概念和推导。
  \item 邱锡鹏的《神经网络与深度学习》适合查阅前馈、卷积、循环、注意力网络及其
    训练方法。
  \item Zhang、Lipton、李沐和 Smola 的 \emph{Dive into Deep Learning}
    适合结合可运行代码理解神经网络的实现、训练与调参过程。
  \item Géron 的 \emph{Hands-On Machine Learning with Scikit-Learn, Keras, and TensorFlow}
    适合查阅从数据准备、模型训练到评估和部署的实践流程。
\end{itemize}

\paragraph*{强化学习}
\begin{itemize}
  \item Sutton 和 Barto 的 \emph{Reinforcement Learning: An Introduction} 适合查阅
    马尔可夫决策过程、价值方法、策略方法与函数逼近。
  \item Szepesvári 的 \emph{Algorithms for Reinforcement Learning} 适合查阅动态规划、
    Monte Carlo、时序差分和函数逼近算法的数学形式与分析。
\end{itemize}

\paragraph*{应用与系统}
\begin{itemize}
  \item Manning、Raghavan 和 Schütze 的 \emph{Introduction to Information Retrieval}
    适合查阅文本表示、索引、检索模型和评价方法。
  \item Huyen 的 \emph{Designing Machine Learning Systems} 从数据、训练、部署与监控
    讨论机器学习系统设计，适合查阅模型进入生产环境后的工程约束。
  \item Jurafsky 和 Martin 的 \emph{Speech and Language Processing}
    适合查阅自然语言与语音处理中的语言模型、序列标注、句法分析和语音识别。
  \item Szeliski 的 \emph{Computer Vision: Algorithms and Applications}
    适合查阅成像、特征、几何、识别与图像合成方法。
  \item Aggarwal 的 \emph{Recommender Systems: The Textbook}
    适合查阅协同过滤、基于内容的推荐、矩阵分解和推荐评价。
  \item Kleppmann 和 Riccomini 的 \emph{Designing Data-Intensive Applications}
    适合查阅存储、复制、分区、事务以及流式与批量处理的系统取舍。
\end{itemize}

上述教材的完整出版信息列于卷末参考文献。
\nocite{mitchell1997,bishop2006,zhou2016machinelearning,li2019statisticallearning,deisenroth2020mathematics}
\nocite{hastie2009,shalevshwartz2014,mohri2018foundations,james2021statistical}
\nocite{pgm-murphy2012,murphy2022probabilistic,murphy2023advanced,pgm-koller2009,rasmussen2006}
\nocite{goodfellow2016,bishop2024deep,qiu2020neural,zhang2023dive,geron2022hands-on}
\nocite{sutton2018,szepesvari2010algorithms}
\nocite{manning2008ir,huyen2022designing,jurafsky2009speech,szeliski2022computer,aggarwal2016recommender,kleppmann2026designing}
